% ============================================================
\chapter{KẾT LUẬN VÀ HƯỚNG PHÁT TRIỂN}
\label{ch:ket_luan}
% ============================================================

\section{Kết luận}

Báo cáo đã nghiên cứu một cách hệ thống bốn chiến lược xử lý mất cân bằng lớp cho bài
toán phân đoạn ngữ nghĩa ảnh UAV rừng, trên cùng kiến trúc UNet++/ResNet50. Các kết
luận chính:

\begin{itemize}[leftmargin=1.4cm]
  \item Bộ dữ liệu rừng UAV có mức \textbf{mất cân bằng lớp rất nghiêm trọng}
  ($2.761\times$), với hai lớp cực hiếm Fence ($0{,}01\%$) và Car ($0{,}04\%$).
  \item \textbf{Trọng số lớp đơn thuần không phải giải pháp tốt}: tuy nâng được lớp
  hiếm nhưng làm giảm hiệu năng lớp phổ biến, dẫn tới mIoU tổng giảm
  ($0{,}8187 \rightarrow 0{,}8142$).
  \item \textbf{Tổ hợp Focal+Dice+CE} cải thiện ổn định ($0{,}8270$), và
  \textbf{phiên bản trọng số thích nghi cho kết quả tốt nhất} ($0{,}8365$,
  $+1{,}78$ pp so với baseline), đặc biệt trên các lớp hiếm (Fence $+10{,}9$ pp,
  Car $+4{,}4$ pp).
  \item Cơ chế \textbf{uncertainty weighting} cân bằng tự động các thành phần loss,
  vừa nâng lớp hiếm vừa bảo toàn lớp phổ biến --- là lựa chọn được khuyến nghị.
\end{itemize}

\section{Hướng phát triển}
\label{sec:huong_morong}

\begin{enumerate}[leftmargin=1.4cm]
  \item \textbf{Hoàn thiện tính chặt chẽ:} đồng nhất số epoch và giao thức đánh giá;
  bổ sung đủ 3 seed cho Exp 4 (Mục~\ref{sec:multiseed}).
  \item \textbf{Độ bền theo điều kiện thời tiết:} hoàn thiện đánh giá trên sunny,
  overcast và dữ liệu trộn (Mục~\ref{sec:weather}).
  \item \textbf{Chuyển miền sang ảnh thực:} đánh giá trên WildUAV
  (Mục~\ref{sec:wilduav}); cân nhắc kỹ thuật thích nghi miền (domain adaptation).
  \item \textbf{So sánh đa kiến trúc:} mở rộng sang CNN hiện đại (ConvNeXt,
  EfficientNetV2), Transformer (SegFormer, SwinV2) và mô hình nhẹ cho thiết bị biên
  (Mục~\ref{sec:multi_arch}).
  \item \textbf{Tăng cường dữ liệu \& hậu xử lý:} copy-paste cho lớp cực hiếm,
  Test-Time Augmentation, làm trơn biên (CRF).
\end{enumerate}
